\chapter{Supplement to Production Study II} 
\label{Appendix B}

\section{Reading material}

\textbf{Story 1}: Banane `banana'\\
\\
Schon seit vielen Jahren arbeitet Susanne im Zoo als Tierpflegerin. Dort ist sie sehr glücklich und kümmert sich liebevoll um alle ihre Schützlinge. Heute hat Susanne \textbf{eine Banane} (1) mitgebracht. Sie hatte eine große Kiste mit viel Obst und Gemüse dabei und besonders die Affen waren ganz aufgeregt. Dann hat Susanne \textbf{die Banane} (2) herausgeholt. Sofort versammelten sich die Affen um sie. \textbf{Die Banane} (3) sah richtig lecker aus und der kleine Affe Tobi wollte sie unbedingt haben. Erst vor ein paar Tagen hatte ein Mädchen mit seinen Eltern vor dem Affengehege gestanden und Tobi beobachtet. Dabei hat das Mädchen \textbf{eine Banane} (4) gegessen. Da konnte Tobi die schöne gelbe Frucht nur ansehen, aber jetzt will er endlich eine probieren! \textbf{Die Banane} (5) wird er bestimmt sehr genießen.\\
\\
\textbf{Story 2}: Gardine `curtain'\\
\\
Sarah hat eine Arbeitsstelle in einer anderen Stadt bekommen und muss jetzt umziehen. Um ihre neue Wohnung einzurichten, war sie im Möbelhaus einkaufen. Dort hat Sarah \textbf{eine Gardine} (1) gefunden. Sie wollte eigentlich erst die Möbel aussuchen und sich dann um die restliche Einrichtung kümmern, sie hat sich aber im Möbelhaus verlaufen. Dabei hat Sarah \textbf{die Gardine} (2) gesehen. Der Preis passt gerade noch in das Budget für die Wohnung. \textbf{Die Gardine} (3) ist hellblau und hat die perfekte Länge für die Fenster im Wohnzimmer. Vorige Woche hatten Sarah und ihre beste Freundin schon zusammen im Internet recherchiert. Dabei hat ihre Freundin \textbf{eine Gardine} (4) vorgeschlagen. Die würde zwar gut ins Schlafzimmer passen, allerdings nicht ins Wohnzimmer. Aber zum Glück ist Sarah ja inzwischen fündig geworden! \textbf{Die Gardine} (5) wird sie noch heute Abend aufhängen.\\
\\
\textbf{Story 3}: Parade `parade'\\
\\
Um seinen Abschluss zu feiern, hat Julian einen Städtetrip unternommen. Er ist mit dem Zug nach Paris gefahren, hat im Hotel eingecheckt und wollte dann die Stadt erkunden. An dem Tag hat Julian \textbf{eine Parade} (1) eingeplant. Er wusste, dass ein Feiertag war und ist sofort in die Innenstadt gefahren. Dort hat Julian \textbf{die Parade} (2) angeschaut. Die Stimmung in der Stadt war super und die Straßen waren voll mit feiernden Menschen. \textbf{Die Parade} (3) war sehr groß und sah professionell aus. Letztes Jahr hat Julian einen Freund in dessen Heimatstadt zum Karneval besucht. Dort hat sein Freund \textbf{eine Parade} (4) organisiert. Die war allerdings sehr klein und ziemlich langweilig. Dieses Mal war Julian aber richtig beeindruckt! \textbf{Die Parade} (5) wird er nie vergessen.\\
\\
\textbf{Story 4}: Garage `garage'\\
\\
Peter wohnt seit einem Jahr in einem neuen Haus, das ihm gut gefällt. Es gibt nur ein Problem: Sein Auto parkt direkt unter einem Baum, wo immer viele Vögel sitzen. Jeden Morgen ist sein Auto dreckig – womit nun allerdings Schluss ist. Heute hat Peter \textbf{eine Garage} (1) gebaut. Er hat lange im Internet nach verschiedenen Bauplänen gesucht. Dann hat Peter \textbf{die Garage} (2) entworfen. Am Wochenende ist er in den Baumarkt gefahren und hat das notwendige Material gekauft. \textbf{Die Garage} (3) ist nämlich aus Metall und relativ groß. Die Idee dazu hat Peter von einem Kollegen bekommen, der vor Kurzem auf einer Architektur-Ausstellung war. Dort hat der Kollege \textbf{eine Garage} (4) gesehen. Peter war sich damals schon sicher, dass er mit all seiner Erfahrung genau so etwas nachbauen könnte. Jetzt hat er es geschafft und ist sehr zufrieden mit dem Ergebnis. \textbf{Die Garage} (5) wird er sicher jeden Tag benutzen.\\
\\
\textbf{Story 5}: Kombüse `galley'\\
\\
Heute war eine Schulklasse zu Besuch im Schifffahrtsmuseum. Ein Mitarbeiter hat den Schülern auf einem Rundgang alles erklärt, was sie wissen wollten. Direkt am Anfang hat die Klasse \textbf{eine Kombüse} (1) besichtigt. In diesem Museum sind nämlich alle Räume auf einem Schiff originalgetreu nachgebildet. Vorsichtig hat die Klasse \textbf{die Kombüse} (2) betreten. Darin waren die Figur eines Kochs, ein Herd, ein Ofen und viele Töpfe und Pfannen zu sehen. \textbf{Die Kombüse} (3) war viel kleiner als die Schüler erwartet hatten. Sie haben auch schon in der Schule einiges über Schiffe gelernt. Dabei hat die Lehrerin \textbf{eine Kombüse} (4) vorgestellt. Die war allerdings Teil eines großen Kriegsschiffs. Ein kleines Schiff wie das im Museum muss den begrenzten Platz strategisch gut ausnutzen und die Kinder wussten, dass viele Dinge an Bord wichtiger waren als das Kochen. \textbf{Die Kombüse} (5) haben sie aber auf jeden Fall interessant gefunden.\\
\\
\textbf{Story 6}: Praline `chocolate candy'\\
\\
Maria hat im Büro eine Kollegin, mit der sie sehr gut befreundet ist. Fast jeden Tag hat sie ein kleines Geschenk für Maria dabei, das sie ihr in der Mittagspause überreicht. Heute hat die Kollegin \textbf{eine Praline} (1) mitgebracht. Maria war noch am Arbeiten, als die Kollegin zu ihr an den Schreibtisch kam. Dann hat die Kollegin \textbf{die Praline} (2) ausgepackt. Sie war von einer edlen Marke und Maria hat sich sehr gefreut. \textbf{Die Praline} (3) war mit Karamell und Nüssen gefüllt. Als Kind war sie mit Süßigkeiten sehr wählerisch, was auch mit einem Erlebnis auf ihrem achten Geburtstag zusammenhängt. Damals hat Maria \textbf{eine Praline} (4) gegessen. Sie hat aber nicht gewusst, dass die manchmal mit Alkohol gefüllt sind. Maria hat sich geekelt und danach sehr lange gar keine Schokolade mehr angerührt. Diese Zeiten sind zum Glück vorbei und heute liebt Maria Schokolade aller Art. \textbf{Die Praline} (5) wird sie gleich als Vorspeise essen.\\
\\
\textbf{Story 7}: Ballade `ballad'\\
\\
Zum Valentinstag ist Max mit seiner Freundin ins Theater gegangen, um sich ein Musical anzusehen. Der Hauptdarsteller war weltbekannt und hatte schon viele Preise gewonnen. Max hat sich schon seit Wochen auf diesen Abend gefreut. Gleich am Anfang hat der Darsteller \textbf{eine Ballade} (1) gesungen. Das Licht wurde gedimmt und der Vorhang ging langsam auf. Dann hat der Darsteller \textbf{die Ballade} (2) begonnen. Er war ein Bariton und hatte eine wunderschöne Stimme. \textbf{Die Ballade} (3) war der Höhepunkt des Musicals. Früher war Max auch in einer Theatergruppe und hat sogar selber schon einmal bei einem Musical mitgemacht. Darin hat Max \textbf{eine Ballade} (4) vorgetragen. Es war eine große Herausforderung und er musste sehr lange üben, um sie zu perfektionieren. Aber der Darsteller in diesem Musical hat scheinbar mühelos gesungen. Das Stück hat Max so gut gefallen, dass er die CD mit der Musik gekauft hat. \textbf{Die Ballade} (5) wird er auf jeden Fall zuhause noch einige Male anhören.

\newpage
\textbf{Story 8}: Limone `lemon'\\
\\
Theo hat eine Nachbarin mit einem sehr großen Garten, die ihre Erzeugnisse regelmäßig an die Nachbarschaft verschenkt. Auch heute stand wieder eine Kiste vor seiner Tür. Unter anderem hat Theo \textbf{eine Limone} (1) bekommen. Der Zeitpunkt war ideal, weil er gerade alle seine Vorräte aufgebraucht hatte. Gut gelaunt hat Theo \textbf{die Limone} (2) herausgeholt. Alles in der Kiste sah sehr appetitlich aus. \textbf{Die Limone} (3) war ganz fest und leuchtend gelb. Theos Freundin geht oft auf dem Markt einkaufen, wo es jetzt im Sommer viele saisonale Angebote gibt. Letzte Woche hat seine Freundin \textbf{eine Limone} (4) gekauft. Die war zwar auch nicht schlecht, aber kein Vergleich zu dem Obst von seiner Nachbarin. Heute will Theo deshalb etwas tun, um sich bei ihr zu bedanken. \textbf{Die Limone} (5) wird er in einem Kuchen verarbeiten und seine Nachbarin zum Kaffee einladen.\\

\largerpage
\textbf{Practice story 1}: Makrele `mackerel'\\
\\
Paul wohnt an der Nordseeküste und an klaren Sommertagen geht er immer fischen. Auch heute ist er mal wieder mit seinem Boot aufs Meer gefahren. Nach einer Weile hat Paul \textbf{eine Makrele} (1) gefangen. Er hat gemerkt, dass sich sein Fischernetz bewegt und hat es eingeholt. Dann hat Paul \textbf{die Makrele} (2) herausgenommen. Sonst hingen nur Plastiktüten und Algen in seinem Netz. \textbf{Die Makrele} (3) war immerhin schön groß und gesund. Er hatte auch schon eine Idee, was er damit machen kann. Letzte Woche war er bei einer Freundin eingeladen. Zum Mittag hat die Freundin \textbf{eine Makrele} (4) gekocht. Die Mahlzeit war sehr lecker und hat ihn schwer beeindruckt. Paul will das Rezept jetzt auch selbst ausprobieren. \textbf{Die Makrele} (5) wird es auf jeden Fall heute zum Abendessen geben.\\
\\
\textbf{Practice story 2}: Rosine `raisin'\\
\\
Emma hatte letzte Woche Besuch von einem Freund, der selbstgebackene Plätzchen mitgebracht hatte. Sie war begeistert und hat sofort ein Plätzchen probiert. Dabei hat Emma \textbf{eine Rosine} (1) geschmeckt. Das war eine unangenehme Überraschung, weil die Stücke in den Plätzchen nach Schokolade aussahen. Dann hat Emma \textbf{die Rosine} (2) ausgespuckt. Es war ihr ein bisschen peinlich, aber ihr Freund war zum Glück gerade abgelenkt. \textbf{Die Rosine} (3) war sehr süß und ein bisschen klebrig. Seit Emma ein Kind war, hat sie eine starke Abneigung gegen getrocknete Früchte. Ihre Mutter hat sie aber immer dazu ermuntert, erstmal alles zu probieren, um sich eine eigene Meinung zu bilden. Also hat Emma \textbf{eine Rosine} (4) gegessen. Danach war sie sehr wütend auf ihre Mutter. Inzwischen kann sie ihren Standpunkt zwar nachvollziehen, aber sie möchte trotzdem nicht alles ausprobieren, was man ihr vorsetzt. \textbf{Die Rosine} (5) wird sie heute sicher nicht mehr essen.

\section{Communicative skill questionnaire}


Subset of questions related to communicative skill from the Autism-Spectrum Quotient questionnaire by \citet{baron-cohenAutismSpectrumQuotientAQ2006}, German translation from \url{https://autismus-kultur.de/autismus-asperger-test-erwachsene/}. Participants respond with “stimme voll zu" (`fully agree'), “stimme eher zu” (`rather agree'), “stimme eher nicht zu” (`rather disagree') or “stimme nicht zu” (`fully disagree').

``Fully disagree” or “rather disagree” responses scored 1 point on the following items: 1, 2, 3, 4, 5, 6. “Fully agree” or “rather agree” responses scored 1 point on the following items: 7, 8, 9, 10.

\begin{enumerate}
\item Andere Menschen sagen mir oft, dass das, was ich gesagt habe, unhöflich war, obwohl ich denke, es sei höflich gewesen.\\
`Other people frequently tell me that what I’ve said is impolite, even though I think it is polite.'
\item Wenn ich rede, ist es für andere nicht immer einfach, zu Wort zu kommen.\\
`When I talk, it isn’t always easy for others to get a word in edgeways.'
\item Ich stelle oft fest, dass ich nicht weiß, wie ich ein Gespräch in Gang halten kann.\\
`I frequently find that I don’t know how to keep a conversation going.'
\item Wenn ich am Telefon mit jemandem spreche, weiß ich nicht genau, wann ich wieder mit Sprechen an der Reihe bin.\\
`When I talk on the phone, I’m not sure when it’s my turn to speak.'
\item Ich verstehe einen Witz oft als Letzte/r.\\
`I am often the last to understand the point of a joke.'
\item Andere Leute sagen mir oft, dass ich ständig über dieselben Dinge spreche.\\
`People often tell me that I keep going on and on about the same thing.'
\item Ich genieße Smalltalk und belangloses Geplauder.\\
`I enjoy social chit-chat.'
\item Es fällt mir leicht, “zwischen den Zeilen zu lesen”, wenn sich jemand mit mir unterhält.\\
`I find it easy to “read between the lines” when someone is talking to me.'
\item Wenn ich mit jemandem rede, merke ich, wenn es ihm/ihr langweilig wird.\\
`I know how to tell if someone listening to me is getting bored.'
\item Smalltalk fällt mir leicht.\\
`I am good at social chit-chat.'
\end{enumerate}
